\documentclass[rjthesis}
\geometry{paperwidth=18.40cm,paperheight=26.00cm,
  left=1.45cm,right=1.45cm,top=1.00cm,bottom=2.20cm}
\usepackage{graphicx}
\graphicspath{{./figures/}}
\usepackage{amsmath}
\usepackage{amsthm}
\newtheorem{definition}{定义}
\newtheorem{property}{性质}
\usepackage{booktabs}
\usepackage[super,square,sort&compress]{natbib}
\usepackage{hyperref}
\hypersetup{colorlinks=true,linkcolor=blue,citecolor=blue}
\usepackage[ruled,vlined,linesnumbered]{algorithm2e}
\usepackage{url}

% ---- Metadata ----
\rjtitle{基于深度学习的文档智能转换方法研究}
\rjauthor{张三, 李四}
\rjinfor{（某某大学计算机科学与技术系，某某 123456）\\
通讯作者: 张三, E-mail: zhangsan@example.edu.cn}

\begin{document}
\rjmakenobibintoc
\rjmaketitle

% ---- Bilingual Abstract ----
\begin{rjabstract}
\AbstractContentZh
本文提出一种基于深度学习的文档智能转换方法。该方法采用双向对齐机制，
在 LaTeX 到 DOCX 转换过程中保持版面语义完整性。实验表明，本文方法在多个
评估指标上优于现有方法。
\end{rjabstract}
\rjkeywords{文档转换; 深度学习; 格式分析; 双向对齐}
\AbstractENContent
{This paper proposes a deep learning based method for intelligent document
conversion. We introduce a bidirectional alignment mechanism to preserve
layout semantics during LaTeX-to-DOCX transformation. Experiments on benchmark
datasets demonstrate superior performance compared to existing methods.}
\rjkeywordsEn{Document Conversion; Deep Learning; Format Analysis; Bidirectional Alignment}

\section{引言}
\label{sec:intro}

文档格式转换是学术出版和协作工作中的常见需求。LaTeX 作为科技文档排版的标准工具，
广泛应用于数学、计算机科学等领域的论文撰写。然而，Microsoft Word 的普及使得
LaTeX 到 DOCX 的格式转换成为必要~\cite{zhang2020deep}。

\section{方法}
\label{sec:method}

\subsection{问题定义}
\label{ssec:problem}

\begin{definition}
\label{def:alignment}
双向对齐（Bidirectional Alignment）指在 LaTeX$\to$DOCX 转换过程中，
源端与目标端语义块之间的双向映射关系，满足对称性：
$\mathcal{A} = \mathcal{A}^{-1}$。
\end{definition}

\subsection{网络结构}
\label{ssec:network}

本文方法采用编码器-解码器架构~\cite{vaswani2017attention}，
如图~\ref{fig:system} 所示。

% ---- Bilingual Figure Caption ----
\begin{figure}[ht]
  \centering
  \includegraphics[width=.7\textwidth]{fig-system-overview.pdf}
  \bicaption{系统总体架构}{System Overall Architecture}
  \label{fig:system}
\end{figure}

% ---- Bilingual Table Caption ----
\begin{table}[ht]
\centering
\bicaption{数据集统计}{Dataset Statistics}
\label{tab:stats}
\begin{tabular}{lccc}
  \toprule
  数据集 & 训练集 & 验证集 & 测试集 \\
  \midrule
  JOS-1000 & 800 & 100 & 100 \\
  IEEE-500  & 400 & 50  & 50  \\
  ACM-300   & 240 & 30  & 30  \\
  \bottomrule
\end{tabular}
\end{table}

\section{实验}
\label{sec:exp}

我们在三个基准数据集上进行了实验~\cite{wang2021table}，
结果如表~\ref{tab:stats} 所示。

\section{结论}
\label{sec:concl}

本文提出了一种基于深度学习的文档智能转换方法~\cite{li2019layout}，
实验结果表明该方法有效。

\bibliographystyle{unsrt}
\bibliography{refs}

\noindent{\hei 附中文参考文献:}
\begin{thebibliography}{99}
\addtolength{\itemsep}{-0.6ex}
\item 张伟, 李明. 基于深度学习的文档格式转换方法[J]. 软件学报, 2020, 31(5): 1234-1245.
\item 王强, 陈华. 表格结构识别研究综述[J]. 计算机学报, 2021, 44(3): 567-578.
\item 李娜, 赵刚. 版面分析在文档理解中的应用[J]. 中文信息学报, 2019, 33(8): 1-10.
\end{thebibliography}

% ---- Author Biography ----
\noindent{\textbf{作者简介:}
\begin{description}[font=\normalfont,labelwidth=2em,leftmargin=2.5em]
\item[{[张三]}] 博士，CCF 专业会员，主要研究方向为文档智能处理与机器学习。
\item[{[李四]}] 硕士，主要研究方向为自然语言处理。
\end{description}}

\end{document}
